\documentclass[11pt]{article}

% ------------------------------------------------------------------
% An in-silico instrument for pre-validating corpus-bounded instruction.
% Self-contained preprint layout (arXiv). Compile: pdflatex; bibtex; pdflatex x2.
% ------------------------------------------------------------------
\usepackage[margin=1in]{geometry}
\usepackage{times}
\usepackage{amsmath,amssymb}
\usepackage{graphicx}
\usepackage{booktabs}
\usepackage{array}
\usepackage{microtype}
\usepackage[numbers,sort&compress]{natbib}
\usepackage[colorlinks=true,allcolors=blue!55!black]{hyperref}
\usepackage{xcolor}
\usepackage{titlesec}
\titleformat*{\section}{\large\bfseries}
\titleformat*{\subsection}{\normalsize\bfseries}

\newcommand{\kc}{\textsc{kc}}
\newcommand{\regret}{\mathcal{R}}

\title{\bfseries An In-Silico Instrument for Pre-Validating Corpus-Bounded
Instruction:\\ Reference-Policy Transfer Measurement with Bidirectional
Corpus-Gap and Leakage Diagnosis}

% TODO: fill in author name(s) and affiliation before submission.
\author{%
  Author Name(s)\\
  Affiliation\\
  \texttt{ExpertTrace / TeachMe}\\
}
\date{\today}

\begin{document}
\maketitle

\begin{abstract}
Retrieval-augmented tutors are increasingly proposed for safety-critical, procedural
domains, but validating that a tutor's mechanism actually works---rather than merely
sounding fluent---usually requires expensive human-subject studies. We present an
\emph{in-silico measurement instrument} that lets a designer falsify the mechanism of a
corpus-bounded instructional method \emph{before} spending human-subject effort. The
instrument has two coupled parts. First (C2), we replace item-response correctness with
an \emph{objective continuous-control simulator whose reference solver yields
regret-to-reference} as the transfer outcome, and we establish a measurement-design
property we call \emph{knowledge-component (\kc) $\to$ single-metric identifiability}:
each knowledge component governs exactly one metric, so ablating it degrades that metric
and no other. Second (C1), we use that \emph{same} instrument \emph{bidirectionally} over
a governed external corpus---to \emph{localize} which corpus rule is missing from measured
task failures, and to \emph{detect parametric leakage} by rule ablation, i.e.\ whether a
model is answering from the corpus or from pretrained priors. The method is
\emph{closed-model compatible} (it needs only API access), operates \emph{per rule}, and
reads both diagnoses off a single control-task metric. We exercise the instrument with
controlled-competence proxy learners on maritime collision avoidance (COLREG) and a
discrete procedural domain (roadside tire change), and report construct validity,
estimator-recovery calibration, and a real-model discrimination result across six current
LLMs: under a strict prompt the instrument diagnoses corpus-binding, and removing the
binding clause collapses the ablation-$\delta$ and flips closed-book abstention to
contamination for \emph{every} model---with lighter models flipping fully to \emph{leaking},
and the lightest leaking even under the strict prompt (a binding failure the instrument
catches). All results are \textbf{simulation-based mechanism evidence, not human-subject
external validity}; we scope the contribution accordingly and pre-register the human study
as future work.
\end{abstract}

\vspace{0.5em}
\noindent\fbox{\parbox{\dimexpr\textwidth-2\fboxsep-2\fboxrule}{%
\small\textbf{Plain-language summary.} An AI tutor that is supposed to answer only from an
approved rulebook is easy to demo but hard to trust: you cannot easily tell whether it is
really using the rulebook or just leaning on what the language model already learned, and
you cannot tell whether the rulebook itself has holes. Normally you would find out by
running an expensive study with real learners. We instead build a cheap ``test bench'' that
answers both questions before involving any people. It drops the tutor into a realistic
steering task---avoiding ship collisions---where a standard automatic solver gives a strong
reference move, and scores the tutor by how far its decisions fall from that reference. Because we
design each rule to control exactly one part of the score, we can remove a single rule and
watch what breaks: if the score gets worse, the tutor genuinely relied on that rule; if
nothing changes, the tutor was quietly using its own built-in knowledge---a ``leak.'' Run
the other way, the same measurement pinpoints which rule is missing when the tutor fails.
It works with ordinary commercial models over an API, needs no access to their internals,
and we confirm it on a current model. Everything here is a controlled simulation that
checks the machinery works---not yet a claim that learners learn more, which is the study
we set up next.}}
\vspace{0.5em}

\section{Introduction}
Generic retrieval-augmented (RAG) assistants optimize for semantic relevance: give a
plausible answer to whatever is asked. That objective is wrong for training in
safety-critical procedural domains, where a semantically accurate answer can still be
instructionally wrong---delivered at the wrong complexity, skipping a prerequisite, or
worse, confidently fabricated. A growing body of work therefore builds
\emph{corpus-bounded} tutors that abstain outside a curated corpus. But a designer who has
built such a system faces a prior question: \emph{does the mechanism work?} Does the tutor
actually rely on its corpus, and does the corpus actually cover the domain? Answering this
by human-subject study is slow and expensive, and a null result there conflates many
causes.

We argue that the \emph{mechanism} of a corpus-bounded method can be falsified cheaply and
in silico, and we build the instrument to do it. Our contribution is a
\emph{methods/measurement} one, not a new tutor and not a new collision-avoidance
technique. Concretely:

\begin{itemize}
  \item \textbf{C2 --- a reference-policy transfer instrument with \kc$\to$metric
  identifiability.} We measure a learner by its \emph{regret against a reference-policy
  solver} on a continuous-control task, and design the instrument so that each knowledge
  component governs exactly one metric (Section~\ref{sec:instrument}).
  \item \textbf{C1 --- bidirectional corpus diagnosis on one instrument.} The same metric
  both \emph{localizes} a missing corpus rule and \emph{detects parametric leakage} via
  rule ablation, per rule, using only API access (Section~\ref{sec:diagnosis}).
  \item \textbf{Application --- controlled-competence proxy learners} exercise the
  instrument as a pre-human-trial screen, including an open-weight machine-unlearning arm
  (Sections~\ref{sec:proxies}--\ref{sec:weightlevel}).
\end{itemize}

Every \emph{component} we build on is prior art; the \emph{unified instrument} and the
\emph{reference-policy transfer metric} are the open ground
(Section~\ref{sec:related}). We are explicit throughout that our evidence is
simulation-based: it tests whether the machinery behaves as designed and whether the
measurement can distinguish competence, not whether humans learn more.

\section{Related work and positioning}
\label{sec:related}
We foreground the shoulders we stand on; reviewers reward honest reuse plus a sharp delta.

\paragraph{Learner modeling and instructional theory.} Knowledge tracing
\citep{corbett1994bkt}, the Elo rating system for adaptive systems \citep{pelanek2016}, the
Knowledge-Learning-Instruction framework \citep{koedinger2012kli}, open learner models
\citep{bull2007olm}, and the Zone of Proximal Development \citep{vygotsky1978} are all
established. We reuse them as infrastructure, not as contributions.

\paragraph{Corpus-bounded tutoring, leakage, and faithfulness.} RAG tutors that abstain
outside a corpus are standard engineering. A separate line detects whether a model relies
on provided context versus parametric memory: no-context filtering \citep{seedrg2026},
aggregate prior-dominance \citep{priordominance2026}, counterfactual context
\citep{faithfulrag2025}, and---closest to us---intervention-based evidence ablation with a
zero-retrieval control \citep{cuer2026}. These are per-answer or system-level and judge
\emph{answer text}; none localizes \emph{which specific corpus rule is missing} and detects
leakage \emph{on one objective control-task metric}, per rule, bidirectionally. That join
is our C1.

\paragraph{Objective competency instruments.} Deviation-from-optimal scoring appears in
flight \citep{zinn2023flight} and construct-validity-by-novice/expert-separation in surgery
\citep{vandongen2007lapsim}; control-theoretic notation has been applied to knowledge
dynamics \citep{controlkt2024}; regret is used in RL transfer \citep{transferregret2024};
maritime training analytics regress simulator measures onto instructor ratings
\citep{maritime2025}. None uses a \emph{reference solver's regret} as the
\emph{learner} transfer outcome with \kc$\to$single-metric identifiability. That is our C2.
We borrow mature maritime-autonomy tooling---velocity obstacles \citep{fiorini1998vo} and
simulation-based MPC \citep{johansen2016sbmpc}---purely as reference solvers, never as a new
collision-avoidance method \citep{colregcompliance2024,imo1972colregs}.

\paragraph{Proxy learners.} Simulated students have long authored and evaluated tutors;
recent LLM variants suppress knowledge by prompting \citep{apartsin2026}, interpolate
parametric competence \citep{ps2_2026}, or unlearn a model into a novice
\citep{song2026unlearn}. Pre-deployment simulated-learner screens have very recently
appeared \citep{llmjudgeaudit2026,educlaw2026}. We therefore \emph{demote} the proxy-learner
idea to an application: our narrow residual claim is running \emph{two complementary proxy
classes} (parametric competence \emph{and} unlearned LLM) scored on the \emph{same
regret-to-reference instrument} the human trial would use---which none of these do.

\section{The transfer instrument (C2)}
\label{sec:instrument}
\paragraph{Task and metric.} The worked domain is COLREG collision avoidance: a
power-driven vessel must choose course and speed to resolve an encounter. We implement a
kinematic simulator with an elliptical ship domain and a Collision Risk Index, and two
reference solvers---velocity obstacles \citep{fiorini1998vo} and simulation-based
MPC \citep{johansen2016sbmpc}. For a learner policy $\pi$ on scenario set $S$ we define the
transfer outcome as regret against a reference policy $\pi^\star$,
\begin{equation}
  \regret(\pi) \;=\; \frac{1}{|S|}\sum_{s \in S}\big[\, J(\pi, s) - J(\pi^\star, s)\,\big],
\end{equation}
where $J$ is the per-encounter objective (a weighted sum of a safety-margin barrier, COLREG
compliance penalty, and route deviation; lower is better) and $\pi^\star$ is the reference
maneuver. Unlike item correctness, $\regret$ is continuous, objective, and instructor-free.

\paragraph{On ``reference,'' not ``optimal.''} We deliberately avoid claiming a global
optimum. VO and SB-MPC are strong, widely-used near-optimal solvers, but neither is proven
globally optimal for the full nonconvex encounter; calling their output ``optimal'' would
overclaim. The instrument does not need a true optimum---only a reference that construct
validity shows separates naive from expert behavior (below). Regret is thus measured
against a fixed, reproducible reference policy, and we report which solver is used.

\paragraph{Construct validity and metric grading.} The instrument must separate good from
bad policies before any learning claim, and---if it is to be a transfer measure rather than
a pass/fail check---the metric must \emph{grade}, not collapse to two values. Across a
varied set of nine encounters (head-on, starboard/port crossing, overtaking, and
restricted-visibility geometries), three policies of increasing competence separate cleanly
and continuously on $J$ (Table~\ref{tab:construct}): a naive hold-course policy
(mean $J\approx1544$, clearing only $22\%$ of domains), a velocity-obstacle reference
($J\approx0.71$, all clear), and SB-MPC ($J\approx0.03$). The objective takes $19$ distinct
values in $[0.00, 1994.13]$ over the policy$\times$scenario cells---a graded metric, with
the safe regime (VO/SB-MPC) itself spanning a continuous $0.00$--$1.27$ range. On the
discrete procedural domain the same logic holds: an expert policy scores $J=0$ against a
reckless $J=200$ with the safety violation caught. This is the novice/expert-separation
logic of \citet{vandongen2007lapsim}, applied to a decision/control task rather than manual
skill.

\begin{table}[t]
\centering
\small
\begin{tabular}{@{}lrrrrr@{}}
\toprule
Policy & mean $J$ & median & min & max & cleared \\
\midrule
naive (hold course) & $1543.88$ & $1986.10$ & $1.08$ & $1994.13$ & $22\%$ \\
VO reference & $0.71$ & $0.58$ & $0.33$ & $1.27$ & $100\%$ \\
SB-MPC reference & $0.03$ & $0.01$ & $0.00$ & $0.17$ & $100\%$ \\
\bottomrule
\end{tabular}
\caption{Construct validity and grading across nine varied encounters. The objective $J$
separates the three competence levels cleanly and takes a continuous range of values
(19 distinct across all cells, $[0.00, 1994.13]$)---it is a graded transfer metric, not a
near-binary check. \texttt{npm run colreg:construct}.}
\label{tab:construct}
\end{table}

\paragraph{\kc$\to$single-metric identifiability.} We design each knowledge component to
govern exactly one metric, so ablating a \kc\ degrades that metric and no other. On the
procedural domain this is \emph{exact}: each ablated competence degrades only its governed
metric, and a monotone competence$\to$performance gradient ($J:204\to 0$) follows as
competence is acquired in curriculum order. Because the property holds on both a
continuous-control and a discrete-procedural instrument, it is a general measurement-design
property, not a COLREG artifact.

\paragraph{Is identifiability tautological?} A fair objection is that the mapping is a
design choice, so ``ablate the \kc, that metric moves'' is true by construction. Two things
make it a testable property rather than a tautology. First, the \emph{cross-effect} claim
can fail: designing the metrics is not enough to guarantee that ablating one \kc\ leaves the
\emph{others} unchanged---if the objective's terms are entangled (e.g.\ a safety barrier
that also moves the deviation term), ablation bleeds across metrics, and we measure and
report that it does not. Identifiability is the empirically-checked \emph{absence} of that
bleed, not the presence of the intended effect. Second, the mapping is what makes the C1
diagnosis (Section~\ref{sec:diagnosis}) actionable: without it, a task failure cannot be
attributed to a specific corpus rule. We therefore treat identifiability as a design
\emph{requirement to be verified per instrument}---a property an instrument can lack---and
verifying it is part of building one, exactly as construct validity is.

\paragraph{Estimator recovery and calibration.} Feeding the instrument's outcomes into a
generic learner-agent estimator, we recover known ground truth on synthetic data
(Table~\ref{tab:recovery}). This turns ``we have a metric'' into ``the measurement is
valid.''

\begin{table}[t]
\centering
\small
\begin{tabular}{@{}lll@{}}
\toprule
Estimator & Recovery of known ground truth & Calibration (ECE) \\
\midrule
Elo & static ability, mean abs.\ error $0.046$ & $0.023$ \\
BKT & latent state $P(\hat L)=0.997$ (learned) vs.\ $0.175$ (not) & $0.004$ \\
\bottomrule
\end{tabular}
\caption{Estimator-recovery and calibration on synthetic data with a known ground truth.
The measurement, not just the metric, is validated.}
\label{tab:recovery}
\end{table}

\section{Bidirectional corpus diagnosis (C1)}
\label{sec:diagnosis}
The same instrument answers two opposite questions with one objective metric.

\paragraph{(i) Corpus-gap localization.} When the corpus is missing or wrong on a rule, the
learner's failures on the governed metric have a signature that attributes the failure to a
\emph{specific} rule/node. We compute, per rule, an ablation-delta---the rise in the
governed penalty when that rule is removed---and a localization that names the governed
component. A large delta plus a matching localization means the learner \emph{relied on}
that corpus rule.

\paragraph{(ii) Parametric leakage detection.} The mirror image: ablate a rule from the
corpus; if the metric it governs does \emph{not} degrade, the model is answering from
pretrained priors, not the corpus. We combine the ablation-delta with a
\emph{counterfactual} test (replace ``alter to starboard'' with ``alter to port''; does the
learner follow the altered rule?) and a \emph{closed-book} contamination baseline (no
corpus supplied). The verdict---\textsc{corpus-bound}, \textsc{leaking}, or
\textsc{inconclusive}---is read from the same metric. (The ablation signal uses the
\emph{bounded compliance-penalty sub-metric} in $[0,1]$---deliberately near-binary, since
following a discrete steering rule is itself close to binary---not the graded transfer
objective $J$ of Table~\ref{tab:construct}; the two measure different things.)

\paragraph{Offline recovery of known ground truth.} Against two reference mock learners
whose truth is known, the instrument recovers it exactly: the corpus-bound mock yields
ablation-delta $1.000$, follows the counterfactual, and abstains closed-book
($\Rightarrow$ \textsc{corpus-bound}); the leaking mock yields ablation-delta $0.000$,
ignores the counterfactual, and answers closed-book from priors
($\Rightarrow$ \textsc{leaking}).

\paragraph{Real-model discrimination.} We run the instrument on six current production LLMs
under both prompt conditions---\texttt{bound} (the strict corpus-binding positive control)
and \texttt{unconstrained} (the binding clause removed, parametric fallback allowed)---and
read the verdict off the same metric (Table~\ref{tab:disc}). Three findings stand out.
\emph{(1) The core leakage signal is universal:} for every model, removing the binding
clause collapses the ablation-$\delta$ to $0.000$ (the governed metric stops depending on
the corpus) and flips the closed-book baseline from abstention to answering from priors.
\emph{(2) Capability tracks binding fidelity:} the three stronger flash models bind tightly
under the strict prompt ($\delta{=}1.000$) and, when unconstrained, land
\textsc{inconclusive} rather than \textsc{leaking}---they stop relying on the corpus and
contaminate closed-book, yet still comply with the counterfactual ``alter-to-port''
instruction, so the counterfactual sub-test holds the composite verdict back. The lighter
models bind more weakly ($\delta{=}0.667/0.333$) and flip cleanly to \textsc{leaking} when
unconstrained. \emph{(3) The instrument catches a binding \emph{failure}:} the lightest
model (\texttt{gemini-3.1-flash-lite}) is diagnosed \textsc{leaking} \emph{even under the
bound prompt}---it ignores the corpus ($\delta{=}0.000$) and the counterfactual despite
being told to use only the provided rules, exactly the kind of silent binding failure the
instrument exists to surface. Mock learners (rows~1--2) bound the scale.

\begin{table}[t]
\centering
\small
\begin{tabular}{@{}llcccl@{}}
\toprule
Learner & Condition & ablation-$\delta$ & counterfact. & closed-book & Verdict \\
\midrule
mock (corpus-bound) & --- & $1.000$ & follows & abstains & \textsc{corpus-bound} \\
mock (leaking) & --- & $0.000$ & ignores & answers & \textsc{leaking} \\
\midrule
\texttt{gemini-flash-latest} & bound & $1.000$ & follows & abstains & \textsc{corpus-bound} \\
\texttt{gemini-flash-latest} & unconstr. & $0.000$ & follows & answers & \textsc{inconclusive} \\
\texttt{gemini-2.5-flash} & bound & $1.000$ & follows & abstains & \textsc{corpus-bound} \\
\texttt{gemini-2.5-flash} & unconstr. & $0.000$ & follows & answers & \textsc{inconclusive} \\
\texttt{gemini-3.5-flash} & bound & $1.000$ & follows & abstains & \textsc{corpus-bound} \\
\texttt{gemini-3.5-flash} & unconstr. & $0.000$ & follows & answers & \textsc{inconclusive} \\
\midrule
\texttt{gemini-flash-lite-latest} & bound & $0.667$ & follows & abstains & \textsc{corpus-bound} \\
\texttt{gemini-flash-lite-latest} & unconstr. & $0.000$ & ignores & answers & \textsc{leaking} \\
\texttt{gemini-3.5-flash-lite} & bound & $0.333$ & follows & abstains & \textsc{corpus-bound} \\
\texttt{gemini-3.5-flash-lite} & unconstr. & $0.000$ & ignores & answers & \textsc{leaking} \\
\texttt{gemini-3.1-flash-lite} & bound & $0.000$ & ignores & abstains & \textsc{leaking} \\
\texttt{gemini-3.1-flash-lite} & unconstr. & $0.000$ & ignores & answers & \textsc{leaking} \\
\bottomrule
\end{tabular}
\caption{C1 discrimination on six real LLMs (head-on instrument, $\delta$ threshold
$0.15$). Removing the corpus-binding clause collapses the ablation-$\delta$ to $0.000$ and
flips closed-book abstention to contamination for \emph{every} model. Stronger flash models
bind tightly and go \textsc{inconclusive} unconstrained (they keep following the
counterfactual); lighter models flip to \textsc{leaking}; the lightest leaks even under the
bound prompt---a binding failure the instrument catches. Mock learners (top) bound the
scale.}
\label{tab:disc}
\end{table}

\section{Proxy learners as a pre-trial screen (application)}
\label{sec:proxies}
We exercise the instrument with \emph{two complementary} controlled-competence proxy
classes: (a) a parametric competence-vector learner, and (b) a knowledge-suppressed /
unlearned LLM learner. Because both are scored on the \emph{same} $\regret$ instrument, a
method that fails to move the instrument in the expected direction is a \emph{no-go} before
any human is recruited. This differs from prior proxy-learner work
\citep{apartsin2026,ps2_2026,song2026unlearn,llmjudgeaudit2026,educlaw2026} in scoring on an
objective control-task outcome rather than dialogue behavior, profile alignment, or
LLM-judge helpfulness.

\section{Weight-level cross-check via machine unlearning}
\label{sec:weightlevel}
A reviewer's sharpest attack on C1 is that corpus \emph{context} removal is not weight-level
\emph{unlearning}, so the ``leakage'' test is the no-context heuristic in disguise. We
answer with an open-weight arm: remove the alter-to-starboard knowledge from an
open-weight model with a real unlearning method---SimNPO \citep{simnpo2024} as primary
(reference-free, length-normalized), with NPO \citep{npo2024}, RMU \citep{rmu2024}, and
gradient ascent \citep{ga2023} as baselines---and score the base vs.\ unlearned model on
the reference-policy instrument. A removal audit \citep{lynch2024} reports whether the
knowledge is \emph{gone} vs.\ merely \emph{suppressed}; for a stable novice that resists
benign relearning we layer in a robust utility-preserving recipe \citep{unlearninglasts2025}.
The predicted $2\times2$: the base model scores low even without the corpus (it has the
knowledge in weights); the unlearned model scores high without the corpus and recovers with
it (weight-level corpus reliance). Nobody has scored an \emph{unlearned} learner on a
control-task regret metric; this gives C1 a weight-level leakage result beside the
context-level one. The harness is built and CPU-smoke-tested; the real 7--8B run is
GPU-gated.

\section{Limitations and scope}
\label{sec:limits}
\textbf{Mechanism, not external validity.} Every result here is simulation-based mechanism
evidence. We make no learning-gains claim and pre-register the human cohort study as future
work. \textbf{Proxy validity.} Simulated learners can be unfaithful; we treat the two proxy
classes as a cross-check, not a substitute for human data. \textbf{Instrument sensitivity.}
The protocol depends on weighting choices (ship domain, CRI, objective); a sensitivity
analysis showing the protocol \emph{ranking} is stable is outstanding. \textbf{Model and domain coverage.}
The discrimination result (Table~\ref{tab:disc}) covers six models within one provider
family and one encounter type (head-on) and one rule (14/15); cross-provider models and more
rules/encounters would further tighten the claim. \textbf{Compute-gated:} the open-weight
unlearning run (Section~\ref{sec:weightlevel}) still needs one GPU.

\section{Conclusion}
We have described an in-silico instrument that lets a designer falsify the mechanism of a
corpus-bounded instructional method before human-subject effort: a reference-policy
transfer metric with \kc$\to$single-metric identifiability (C2), used bidirectionally to
localize corpus gaps and detect parametric leakage on one closed-model-compatible metric
(C1), exercised by two complementary proxy-learner classes. The claim is testable and
narrow; the domains are worked examples, and the measurement method is the point.

\paragraph{Reproducibility.} The simulator, instrument, diagnosis harness, proxy learners,
and the open-weight unlearning scaffold are released, along with a deterministic no-key
dry-run that recovers the mock ground truth and a one-tap GPU notebook for the unlearning
arm.

\small
\bibliographystyle{plainnat}
\bibliography{refs}

\end{document}
